\subsection{Exact ordinary-nullification tests for the reversible closures}
\label{sec:A34-reduction}

The distinction between ordinary acyclicity and reversible asphericity
is essential in the two questions above.  We reprove the bounded
suspension comparison used in the supplied REVLOC D.2.7 and D.2.9.
The comparison is an equality of classes; it is not an assumption that
ordinary nullification already has reverse excision.

\begin{claim}{A34.1}{Proved}{A bounded suspension computes the reversible closure}\label{A341}
Let $A$ be a connected pointed space with $p$-primary reduced integral
homology.  Suppose that for an integer $d\geq1$, ordinary $A$-nullification
kills every $d$-connected generalized Eilenberg--Mac Lane space with
$p$-primary homotopy groups.  Write $W_A=\min\operatorname{Rev}(A\to *)$.
Then, without cardinality or torsion-exponent restrictions,
\[
 \mathcal A_{W_A}
   =\{X:P_A(\Sigma^dX)\simeq *\}.
 \tag{A34.1.1}
\]
In particular $P_A(\Sigma^dX)\not\simeq *$ excludes $X$ from the
\emph{reversible} aspherical class.  Failure of $P_AX\simeq *$ alone
does not establish that exclusion.
\end{claim}
\begin{proof}
Put $\mathcal C_A=\{Y:P_AY\simeq *\}$ and
$\mathcal N_d=\{X:\Sigma^dX\in\mathcal C_A\}$.
Existence and the mapping-space universal property of ordinary
nullification are \citet[Theorem 1.1]{ABousfield1996}.
Equivalently, $Y\in\mathcal C_A$ means that
$\operatorname{Map}_*(Y,Z)$ is contractible for every $A$-null $Z$.
This characterization proves closure under weak equivalences, retracts,
and pointed homotopy colimits.  Applied to the mapping-space fiber
sequence of a cofiber sequence $U\to V\to C$, it also gives the two
implications $(U,V)\Rightarrow C$ and $(U,C)\Rightarrow V$.
Suspension preserves pointed homotopy colimits and cofibers, so these
properties hold for $\mathcal N_d$ too.

First observe that every member of $\mathcal C_A$ has $p$-primary
reduced integral homology.  The collapse of $A$ is an
$H\mathbb Z[1/p]$-equivalence, so every space local for that homology
theory is $A$-null.  The mapping-space test just given therefore makes
$Y\to *$ an $H\mathbb Z[1/p]$-equivalence whenever
$Y\in\mathcal C_A$.  Flatness of $\mathbb Z[1/p]$ says precisely that
every element of $\widetilde H_*(Y;\mathbb Z)$ is killed by a power of
$p$.  No uniform bound on those powers is used.

The nonformal point is closure under desuspension.  Suppose that
$\Sigma X\in\mathcal N_d$, or equivalently that
$\Sigma^{d+1}X\in\mathcal C_A$.  The preceding observation and the
suspension homology isomorphism give $p$-primary reduced integral
homology for $X$.  In particular $X$ is connected, because its reduced
$H_0$ is free abelian.  Hence $\Sigma^dX$ is $d$-connected.
The Dold--Thom description makes
$SP^\infty\Sigma^dX$ a $d$-connected GEM with $p$-primary homotopy
groups, so the hypothesis puts it in $\mathcal C_A$.

The exact classical identity is
\citet[Theorem 3.3]{ABousfield1996}: for connected $U$ and $k\geq1$,
\[
 \langle\Sigma U\rangle
   =\langle\Sigma^kU\vee SP^\infty\Sigma U\rangle.
 \tag{A34.1.2}
\]
These are ordinary nullity classes.  Equality means that the two
spaces kill each other and give equivalent nullification functors;
it is not an equality of stable localizing subcategories.  Apply it to
$U=\Sigma^{d-1}X$ and $k=2$.  Both summands on the right are
$A$-acyclic.  Their wedge therefore is, and transitivity of ordinary
acyclicity gives $\Sigma^dX\in\mathcal C_A$.  This proves the required
implication $\Sigma X\in\mathcal N_d\Rightarrow X\in\mathcal N_d$.
The converse follows from closure under suspension.

If $V,C\in\mathcal N_d$ in $U\to V\to C$, rotate to
$V\to C\to\Sigma U$.  The first cofiber implication gives
$\Sigma U\in\mathcal N_d$, and desuspension gives
$U\in\mathcal N_d$.  Thus $\mathcal N_d$ has all three cofiber
implications and all the closure properties in A.7.  That construction
therefore gives a reversible fundamental localizer with aspherical
class $\mathcal N_d$.  It contains $A\to *$, since all positive
suspensions of $A$ are ordinarily $A$-acyclic.  Minimality yields
$\mathcal A_{W_A}\subseteq\mathcal N_d$.

For the other inclusion, every reversible fundamental localizer
containing $A\to *$ inverts ordinary $A$-nullification units.  This
can be seen directly from the cone-attachment construction: its new
cofibers are suspensions of $A$, and the trivial cofibrations and
homotopy colimits used in the construction are already inverted.
Consequently its aspherical class contains $\mathcal C_A$.
For $X\in\mathcal N_d$, the space $\Sigma^dX$ lies in
$\mathcal C_A$; applying reversible desuspension $d$ times puts $X$
in that class as well.  Apply this to $W_A$ to obtain the reverse
inclusion in \textup{(A34.1.1)}.
\end{proof}

\begin{claim}{A34.2}{Proved}{The suspension bounds in the two open problems}\label{A342}
For a prime $p$, put $B=K(\mathbb F_p,2)$ and $A_r=Q\Sigma^rB$,
where $Q=\Omega^\infty\Sigma^\infty$ and $r\geq0$.  The following
are exact equalities of pointed aspherical classes:
\[
\begin{aligned}
 \mathcal A_{\min\operatorname{Rev}(B)}
   &=\{X:P_B(\Sigma X)\simeq *\},\\
 \mathcal A_{\min\operatorname{Rev}(BC_p\vee V)}
   &=\{X:P_{BC_p\vee V}(\Sigma X)\simeq *\},\\
 \mathcal A_{\min\operatorname{Rev}(A_r)}
   &=\{X:P_{A_r}(\Sigma^{r+1}X)\simeq *\}.
\end{aligned}
\tag{A34.2.1}
\]
In the middle equality $V$ is any connected space with $p$-primary
reduced integral homology, including $V=*$.  The first two equalities
are the instances needed from REVLOC D.2.7(a,b); the third gives the
test in REVLOC D.2.9(b).
\end{claim}
\begin{proof}
The exact primary input is
\citet[Lemma 8.13]{ABousfield1997}.  If $n\geq2$ and $G$ is a
nonzero $p$-torsion abelian group, $K(G,n)$ kills every
$(n-1)$-connected generalized polyGEM with $p$-torsion homotopy
groups, with the extra requirement
$\pi_nY/p\pi_nY=0$ when $G/pG=0$.  A GEM is a generalized polyGEM.
For $G=\mathbb F_p$ the exceptional requirement is absent.
Thus $B$ kills every simply connected $p$-primary GEM, and A34.1
applies with $d=1$.

The group space $BC_p$ kills its delooping $K(\mathbb F_p,2)$.
Indeed the usual simplicial bar model of the delooping has finite
powers of $BC_p$ as its levels.  Ordinary nullification preserves
finite products, as stated in the proof of
\citet[Theorem 1.7]{ABousfield1996}; its acyclic class is closed
under realization.  These facts prove the assertion about the
delooping.  Transitivity and Lemma 8.13 then show that $BC_p$ kills
every simply connected $p$-primary GEM.  A wedge $BC_p\vee V$ kills
its $BC_p$ retract and hence the same GEMs.  Its reduced homology is
still $p$-primary, so A34.1 again applies with $d=1$.

For the last equality, $A_r$ is $(r+1)$-connected and
\[
 H_{r+2}(A_r;\mathbb Z)=\mathbb F_p.
\]
This follows from the stable Hurewicz theorem applied to
$\Sigma^r\Sigma^\infty B$, whose first nonzero homotopy group is
in degree $r+2$ and is $\mathbb F_p$.  All reduced homology of $A_r$
is $p$-primary.  One way to see the latter assertion is the Snaith
splitting used in A.6: after inverting $p$, every extended-power
summand of $\Sigma^r\Sigma^\infty B$ vanishes.
By \citet[Corollary 3.2(ii)]{ABousfield1996}, $A_r$ kills
$SP^\infty A_r$.  The Dold--Thom GEM has
$K(\mathbb F_p,r+2)$ as a homotopy retract, using the displayed
bottom homology group.  Hence $A_r$ kills that Eilenberg--Mac Lane
space, which in turn kills every $(r+1)$-connected $p$-primary GEM
by Lemma 8.13.  Apply A34.1 with $d=r+1$.
\end{proof}

\paragraph{The resulting exact gates.}
For the example $L_p=\operatorname{loc}(\Sigma^\infty B)$ of A3.1,
a space $X$ with
\[
 \Sigma^\infty X\in L_p,
 \qquad P_B(\Sigma X)\not\simeq *
 \tag{A34.G1}
\]
would give an actual functor outside the least reversible localizer
and inside its $RT$ closure, namely
$\Delta/\operatorname{Sing}X\to e$.  For the descending candidates
$W_r=\min\operatorname{Rev}(A_r\to *)$, whose common stable part is
$L_p$, the corresponding test on the original generator is exactly
\[
 B\notin\mathcal A_{W_r}
 \quad\Longleftrightarrow\quad
 P_{A_r}(\Sigma^{r+1}B)\not\simeq *.
 \tag{A34.G2}
\]
The statement about the common stable part follows from the Snaith
splitting: its linear summand is
$\Sigma^r\Sigma^\infty B$ and every higher extended power belongs
to the same localizing tensor ideal.  This stable observation does
not evaluate the unstable nullification in \textup{(A34.G2)}.

For A.4 at $p=2$, a space $X\in\Nul_2$ with
$\Sigma^\infty X\in\operatorname{loc}(\mathbb S/2)$ lies outside
$\mathcal A_{\min\operatorname{Rev}(BC_2)}$ if and only if
\[
                 P_{BC_2}(\Sigma X)\not\simeq *.
 \tag{A34.G3}
\]
The one initial suspension in \textup{(A34.G1)} and
\textup{(A34.G3)}, and the $r+1$ suspensions in
\textup{(A34.G2)}, are part of the proved comparisons.  None can be
dropped merely because ordinary nullification is available.
No space satisfying any of these separating tests has been established
in this run; the formulas sharpen the open questions rather than solve
them.

\begin{claim}{A34.3}{Proved}{Smashing the classifying-space example with the higher generator collapses}\label{A343}
The following is an independent proof of the corresponding case of
the supplied REVLOC D.2.37, using a stabilization-cofiber argument.
Let $C=BC_p$, $B=K(\mathbb F_p,2)$, and $A=C\wedge B$, with
derived pointed smash products.  Then
\[
 \operatorname{loc}(\Sigma^\infty A)
   =\operatorname{loc}(\Sigma^\infty B)=L_p,
 \qquad
 P_A(\Sigma^2B)\simeq *,
 \tag{A34.3.1}
\]
and consequently
\[
 \min\operatorname{Rev}(A\to *)
   =\min\operatorname{Rev}(B\to *).
 \tag{A34.3.2}
\]
In particular, if $M=M(\mathbb F_p,1)$ is the cofiber of the
degree-$p$ map of $S^1$, then $M\wedge B$ belongs to this reversible
aspherical class.  Thus smashing the known $BC_p$ example and its
Moore-space witness with $B$ does not yield a new separation for A.3.
This is an exclusion of that specified candidate, not a decision of
the existential question A.3.
\end{claim}
\begin{proof}
We first verify the stable comparison.  Put $D=\Sigma^\infty B$.
\citet[Theorem~1.5]{ALevy} gives, in particular,
$\operatorname{loc}(\Sigma^\infty C)
 =\operatorname{loc}(\mathbb S/p)$.
Smashing the constructions for the two inclusions with $D$ gives
\[
 \operatorname{loc}(\Sigma^\infty A)
   =\operatorname{loc}((\mathbb S/p)\wedge D).
\]
These constructions use shifts, cofibers, coproducts, and retracts,
all preserved by the derived smash product.  Since $D$ is connective
and has $p$-primary integral homology, $D[1/p]$ is connective and
integral-homology acyclic, hence zero by stable Hurewicz.  Every
$D/p^m$ is a finite extension of copies of $D/p$ by the octahedral
axiom.  The maps from the cofiber of $p^m:D\to D$ to the cofiber of
$p^{m+1}:D\to D$, induced by $1_D$ on the source and $p$ on the
target, give a cofiber sequence
\[
 D\longrightarrow D[1/p]
   \longrightarrow\operatorname*{hocolim}_{m\geq1}D/p^m.
\]
Thus $\Sigma D$ lies in $\operatorname{loc}(D/p)$, proving the
reverse inclusion needed for the first equality in
\textup{(A34.3.1)}.  This is an equality of localizing subcategories,
not merely an agreement of Morava-homology vanishing patterns.

We next control an unstable stabilization cofiber.  We use the
published cellular inequalities of
\citet[Propositions~3.1(1),(2) and~3.4(11)]{A4CFFS2015}.
Here ``$U$ is $V$-cellular'' means that $U$ belongs to the smallest
class containing $V$ and closed under pointed homotopy colimits and
weak equivalences; such a class is also closed under homotopy
retracts.  Since $\Omega B\simeq C$, Proposition~3.1(2) says that
$B$ is $\Sigma C$-cellular.  Smashing with $B$ and using
Proposition~3.1(1) shows that
\[
             B\wedge B\text{ is }\Sigma A\text{-cellular}.
 \tag{A34.3.3}
\]
For every connected pointed $Z$ and every $n\geq1$,
Proposition~3.4(11) says that the homotopy cofiber of the adjunction
unit
\[
          e_{n,Z}:Z\longrightarrow\Omega^n\Sigma^n Z
\]
is $Z\wedge Z$-cellular.  The units are compatible with the usual
stabilization maps, and the sequential homotopy colimit of their
targets is $QZ=\Omega^\infty\Sigma^\infty Z$.  Homotopy cofibers
commute with this homotopy colimit.  Consequently
\[
       \operatorname{Cof}(Z\longrightarrow QZ)
          \text{ is }Z\wedge Z\text{-cellular}.
 \tag{A34.3.4}
\]
This uses a direct homotopy colimit, not preservation of acyclicity
under a Postnikov or Goodwillie inverse limit.  Taking $Z=\Sigma B$
and applying \textup{(A34.3.3)} shows that
\[
 K:=\operatorname{Cof}(\Sigma B\longrightarrow Q\Sigma B)
       \text{ is }A\text{-cellular},
 \tag{A34.3.5}
\]
since $(\Sigma B)\wedge(\Sigma B)=\Sigma^2(B\wedge B)$ is
$\Sigma^3A$-cellular and hence $A$-cellular.

For completeness, we verify the spectral-nullification input in
exactly the connectivity range used here.  Write $F=\Sigma^\infty A$.
The space $A$ is $2$-connected, all its reduced integral homology is
$p$-primary, and the reduced K\"unneth formula gives
\[
                  H_3(A;\mathbb Z)=\mathbb F_p.
\]
Thus $F$ is $3$-connective and satisfies the $3$-supported
$p$-torsion condition in
\citet[Theorem~2.6]{ABousfield1996}.  Let
$E\in\operatorname{loc}(F)\cap\mathrm{Sp}_{\geq3}$ and put
\[
 F_i=\Sigma^{-i}F,
 \qquad Y_i=P^{\mathrm{Sp}}_{F_i}E,
 \qquad i\geq0.
\]
Spectral nullification here uses the mapping space, or equivalently
the nonnegative suspensions of the indicated generator.  The cited
theorem supplies natural maps $Y_i\to Y_{i+1}$ with
\[
          \operatorname{Fib}(Y_i\to Y_{i+1})
                 \simeq\Sigma^{2-i}HG_i
 \tag{A34.3.6}
\]
for $p$-primary abelian groups $G_i$.  Exactness of filtered colimits
of spectra and extension closure of the upper homotopy bound imply,
for $Y_\infty=\operatorname*{colim}_iY_i$,
\[
        \operatorname{Fib}(Y_i\to Y_\infty)
                    \in\mathrm{Sp}_{\leq2-i}.
\]
To control components as well as positive homotopy groups, use the
map from $Y_{i+1}$.  Its fiber has upper bound $1-i$, and its
suspended fiber has upper bound $2-i$, whereas
$F_i\in\mathrm{Sp}_{\geq3-i}$.  The orthogonality of the standard
$t$-structure therefore gives an equivalence
\[
 \operatorname{Map}_{\mathrm{Sp}}(F_i,Y_{i+1})
       \simeq\operatorname{Map}_{\mathrm{Sp}}(F_i,Y_\infty).
\]
The left side is contractible because $Y_{i+1}$ is $F_{i+1}$-null
and $F_i=\Sigma F_{i+1}$.  Hence $Y_\infty$ is right orthogonal to
every integer shift of $F$.  The cone-attachment construction of
nullification puts each $Y_i$, and therefore $Y_\infty$, in
$\operatorname{loc}(F)$, since $E$ lies there.  Orthogonality to its
own identity gives $Y_\infty=0$.  The fiber bound at $i=0$ now gives
$Y_0\in\mathrm{Sp}_{\leq2}$.  Nullification by the $3$-connective
$F$ preserves $3$-connectivity: all the attached cones are on
nonnegative suspensions of $F$.  Thus $Y_0=0$.  We have proved
\[
 E\in\operatorname{loc}(F)\cap\mathrm{Sp}_{\geq3}
       \quad\Longrightarrow\quad P_F^{\mathrm{Sp}}E=0.
 \tag{A34.3.7}
\]
No compactness of $F$ or interchange of
$\operatorname{Map}(F,-)$ with a filtered colimit was used.

Now take $E=\Sigma^\infty\Sigma B$.  The stable comparison places
it in $\operatorname{loc}(F)$, and it is $3$-connective.
\citet[Theorem~2.10]{ABousfield1996} identifies
\[
 P_A(Q\Sigma B)
  \simeq\Omega^\infty P_F^{\mathrm{Sp}}
                      (\Sigma^\infty\Sigma B)
  \simeq *.
 \tag{A34.3.8}
\]
The cofiber $K$ in \textup{(A34.3.5)} is also ordinarily
$A$-acyclic, since cellular spaces are acyclic.  Rotate the cofiber
sequence to obtain
\[
                     Q\Sigma B\longrightarrow K
                                \longrightarrow\Sigma^2B.
\]
Ordinary $A$-acyclic spaces are closed under homotopy cofibers.
The first two terms are acyclic, so the third is, proving the
second assertion in \textup{(A34.3.1)}.

Every reversible fundamental localizer containing $A\to *$
contains ordinary $A$-acyclic spaces.  Reversible suspension
reflection therefore makes $B$ aspherical for
$\min\operatorname{Rev}(A\to *)$.  Conversely $A=C\wedge B$
belongs to the reversible class generated by $B$, by pointed
homotopy-colimit and smash closure.  Minimality gives
\textup{(A34.3.2)}, and the same smash closure puts $M\wedge B$
in this class.  The argument does not transport the old Moore
separator through smash product; it proves that its smash is
aspherical in the new candidate.
\end{proof}

\subsection{The degree-two Pr\"ufer comparison is strictly smaller stably}
\label{sec:A34-prufer}

\begin{claim}{A34.4}{Proved}{Cyclotomic completion separates the two degree-two generators}\label{A344}
For a prime $p$, put
\[
 B=K(\F_p,2),\qquad D=K(C_{p^\infty},2),\qquad
 J_p=\loc(\Sigma^\infty D),\quad
 L_p=\loc(\Sigma^\infty B).
\]
There is a spectrum $E_p$, specified in the proof, such that
\[
 E_p\wedge\Sigma^\infty D=0,
 \qquad E_p\wedge\Sigma^\infty B\ne0.
 \tag{A34.4.1}
\]
Consequently
\[
 0\ne J_p\subsetneq L_p,\qquad
 J_p\cap\Ho(\Sp)^\omega=0,
 \qquad
 \min\operatorname{Rev}(D\to*)
 \subsetneq\min\operatorname{Rev}(B\to*).
 \tag{A34.4.2}
\]
Moreover $P_D\Sigma^kB\not\simeq *$ for every $k\geq0$.
These comparisons do not furnish an A.3 witness for the generator $D$:
the displayed separating space $B$ is outside its stable part $J_p$.
\end{claim}
\begin{proof}
All tensor products without a subscript in this proof are ordinary
smash products of spectra.  Fix a telescope $T=T(2)$ at the chosen
prime and write
\[
 \mathcal C=\Sp_{T},\quad L_T:\Sp\longrightarrow\mathcal C,
 \quad U\otimes_{\mathcal C}V=L_T(U\wedge V),\quad
 \mathbf1[X]=L_T\Sigma^\infty_+X.
\]
This category is infinitely semiadditive and has semiadditive height
$2$: the exact primary inputs are
\citet[arXiv v3, Theorem 5.3.1]{A3CSYAmbidexterity2022}
and \citet[arXiv v2, Theorem 4.4.5]{A3CSYHeight2020}.
Let
\[
 R_r=\mathbf1[\omega^{(2)}_{p^r}],\qquad
 R_\infty=\mathop{\operatorname{colim}}_{r\in\mathbb N}^{\mathcal C}R_r
\]
be the height-two cyclotomic algebras of
\citet[Definitions 4.7 and 4.10]{A3CSYCyclotomic2024}.

We first specify a nonzero object annihilated by $R_\infty$.
Use the form of $BP\langle1\rangle$, the $\mathbb Z$-action,
and the completion conventions of
\citet[Convention 5.1 and Theorem 5.4]{A3BHLS2023}.
There is the coassembly map
\[
 c:N=L_TK\bigl(BP\langle1\rangle^{h\mathbb Z}\bigr)
 \longrightarrow
 N'=\bigl(L_TK(BP\langle1\rangle)\bigr)^{h\mathbb Z}.
 \tag{A34.4.3}
\]
By \citet[Theorem 6.25, with $n=1$ and $k=0$]{A3BHLS2023},
this map is not an equivalence and is the cyclotomic-completion
map.  Section 1.1 of the same primary source identifies that
completion as localization in $\mathcal C$ with respect to
$R_\infty$.  Thus the actual fiber
\[
 F=\operatorname{fib}(c)\in\mathcal C
 \quad\text{satisfies}\quad
 F\ne0,\qquad F\otimes_{\mathcal C}R_\infty=0.
 \tag{A34.4.4}
\]
This uses the stronger cyclotomic conclusion of that theorem,
not merely the failure of $K(2)$-locality.

The proof of \citet[Proposition 4.11, pp.~31--32 in arXiv v2]{A3CSYCyclotomic2024}
gives an idempotent splitting for the actual map induced by
$p:D\to D$.  In underlying objects of $\mathcal C$ it identifies
that map with a projection
\[
 \mathbf1[D]\simeq\mathbf1[D]\oplus R_\infty
 \longrightarrow\mathbf1[D].
 \tag{A34.4.5}
\]
The source proof obtains this first on $B^2C_{p^\infty}$ by a
filtered colimit; its further comparison with $B^3\mathbb Z_p$
is not needed here.  Tensoring \textup{(A34.4.5)} with $F$ and
using \textup{(A34.4.4)} shows that
\[
 L_T(F\wedge\Sigma^\infty_+p)
 :L_T(F\wedge\Sigma^\infty_+D)
 \longrightarrow L_T(F\wedge\Sigma^\infty_+D)
\]
is an equivalence.

Now set $E_p=T\wedge F$ in ordinary spectra; no ring structure on
$E_p$ is asserted or used.  A map becomes an
equivalence under $L_T$ precisely when smashing it with $T$
is an equivalence.  Hence $E_p\wedge\Sigma^\infty_+p$ is an
ordinary equivalence.  On the other hand
\[
 \operatorname*{hocolim}
 (D\xrightarrow{p}D\xrightarrow{p}D\longrightarrow\cdots)
 \simeq *:
\]
its only possible nonzero homotopy group is the direct limit of
$C_{p^\infty}$ under multiplication by $p$, and every element
becomes zero.  Smashing the suspension spectrum of this telescope
with $E_p$ commutes with its homotopy colimit.  Since all its
transition maps are equivalences, the canonical augmentation
$E_p\wedge\Sigma^\infty_+D\to E_p$ is an equivalence.
Its fiber is $E_p\wedge\Sigma^\infty D$, proving the first
assertion of \textup{(A34.4.1)}.  This passage through $T\wedge(-)$
is essential: no unlocalized smash product was identified with
the tensor product in $\mathcal C$.

For the other assertion,
\citet[Proposition 4.5 with $r=1$, and Proposition 4.9(2)]{A3CSYCyclotomic2024}
identify $L_T\Sigma^\infty B$ with $R_1$ and show that $R_1$
is faithful in $\mathcal C$.  Thus
\[
 L_T(F\wedge\Sigma^\infty B)
 \simeq F\otimes_{\mathcal C}R_1\ne0.
\]
By the defining test for $T$-acyclicity this says exactly that
$T\wedge F\wedge\Sigma^\infty B\ne0$, as required.

By \citet[Lemma 8.13]{ABousfield1997}, $B$ ordinarily kills $D$,
so $J_p\subseteq L_p$.  The kernel of the exact
colimit-preserving functor $E_p\wedge(-)$ contains $J_p$ but
does not contain $\Sigma^\infty B$, proving strictness.
Also $H_2(D;\mathbb Z)=C_{p^\infty}\ne0$, and the assertion
about finite spectra follows from $J_p\subseteq L_p$ and
Claim~\ref{A31}.  The minimal-construction argument of A.7 gives
$T(\min\operatorname{Rev}(D\to*))=J_p$ and similarly gives
$L_p$ for $B$, so the inclusion of reversible localizers is strict.

Finally, every ordinarily $D$-acyclic space is $E_p$-acyclic,
since $D\to *$ is an $E_p$-equivalence and every
$E_p$-local space is $D$-null.  The nonzero spectrum
$E_p\wedge\Sigma^{\infty+k}B$ therefore prevents
$P_D\Sigma^kB$ from being contractible for every $k\geq0$.
\end{proof}

\paragraph{The remaining A.3 test for this generator.}
Lemma 8.13 of \citet{ABousfield1997} says that $D$ kills every
$2$-connected $p$-primary GEM; its divisible-group condition
holds because the degree-two homotopy group is zero.
Claim~\ref{A341} therefore gives
\[
 \mathcal A_{\min\operatorname{Rev}(D\to*)}
   =\{X:P_D\Sigma^2X\simeq *\}.
 \tag{A34.4.6}
\]
An A.3 witness for $D$ would require both
$\Sigma^\infty X\in J_p$ and $P_D\Sigma^2X\not\simeq *$.
Claim~\ref{A344} computes the second condition for $X=B$ and
refutes its first condition.  It therefore removes that proposed
same-stable-part comparison and leaves the A.3 existence question open.
The deduction above is recorded as a proved comparison in this run,
without a priority claim.  Its inputs are the cited primary theorems;
a separate prior-art search is recorded in the source and attempt ledgers.

\subsection{A universal reduction of A.3 to one prime}
\label{sec:A34-prime-reduction}

\begin{claim}{A34.5}{Proved}{Arithmetic components detect reversible asphericity and failure of $RT$}\label{A345}
For a prime $p$ put $I_p=\loc(S/p)$ and
\[
 \Gamma_pE=\operatorname{fib}(E\longrightarrow E[1/p]),\qquad
 \Gamma_{\mathrm{tor}}E=\operatorname{fib}(E\longrightarrow E_{\mathbb Q}).
\]
Let $W$ be any reversible fundamental localizer, with no accessibility
hypothesis, and define
\[
 W_p=W\cap R(I_p),\qquad
 W_{\mathbb Q}=W\cap R(\loc(H\mathbb Q)).
\]
The complete membership test for $W_p$ is that $g\in W$ and $Bg$
is an $H\mathbb Z[1/p]$-homology equivalence.  These are reversible
fundamental localizers, and
\[
 \begin{split}
 T(W_p)&=T(W)\cap I_p=\Gamma_pT(W),\\
 T(W_{\mathbb Q})&=T(W)\cap\loc(H\mathbb Q),\\
 W&=\bigvee\bigl(\{W_{\mathbb Q}\}\cup\{W_p:p\text{ prime}\}\bigr).
 \end{split}
 \tag{A34.5.1}
\]
Here $\Gamma_pT(W)$ denotes the class of images of its objects;
it is a localizing subcategory.  The rational component is either
$W_{\mathbb Z}$ or $R(\loc(H\mathbb Q))$.  In particular
\[
 W=R(T(W))
 \quad\Longleftrightarrow\quad
 W_p=R(T(W_p))\ \text{for every prime }p.
 \tag{A34.5.2}
\]
Consequently A.3 is equivalent to its restriction to localizers
$V$ satisfying $T(V)\subseteq I_p$ for one prime $p$.
This is a reduction of the existential question, not a construction
of a counterexample.
\end{claim}
\begin{proof}
Write $\mathcal A_W$ for the pointed aspherical class of $W$.
We establish the space-level comparison explicitly; arbitrary fibers
of maps between aspherical spaces are not used.

\paragraph{Localizations by actual telescope maps.}
For a simply connected suspension $Z$ and $m\geq1$, let
$[m]_Z:Z\to\bigvee_{j=1}^m Z\to Z$ be pinch followed by fold.
Each projection of the pinch is the identity, and $[m]_Z$ acts
as multiplication by $m$ on reduced integral homology.
For a set $S$ of primes choose a sequence of integers with prime
factors in $S$, each member of $S$ occurring infinitely often.
The telescope has homology
$\widetilde H_*(Z;\mathbb Z)\otimes\mathbb Z[S^{-1}]$ and is
simply connected.  The homological characterization and uniqueness
in \citet[Theorems 6.1.1--6.1.2]{A3MayPonto2012} identify it,
with its canonical map, as $Z[S^{-1}]$.
This independently proves the needed suspension case of their
Proposition 6.6.4 by homology, without its surrounding assertion
about multiplication on every unstable homotopy group.

Hence $Z\in\mathcal A_W$ implies
$Z[S^{-1}]\in\mathcal A_W$, by closure under pointed homotopy
colimits.  Suspension spectra turn the copower into the actual
scalar $m$, so, compatibly with the localization units,
\[
 \Sigma^\infty Z[S^{-1}]
       \simeq(\Sigma^\infty Z)[S^{-1}].
 \tag{A34.5.3}
\]
This applies to $\mathbb Z[1/p]$, $\mathbb Z_{(p)}$, and
$\mathbb Q$.  No claim about scalar multiplication on all unstable
homotopy groups enters this construction.

\paragraph{The relevant fibers are also cofibers.}
For any pointed space $X$ put $Z=\Sigma^3X$.  It is
$2$-connected, irrespective of the components of $X$.  Define
\[
 G_p(Z)=\operatorname{fib}(Z\longrightarrow Z[1/p]),\qquad
 F(Z)=\operatorname{fib}(Z\longrightarrow Z_{\mathbb Q}).
\]
The homotopy exact sequence and localization on homotopy groups
show that these fibers are simply connected, that $G_p(Z)$ has
$p$-primary homotopy groups, and that $F(Z)$ has torsion homotopy
groups.  Their reduced integral homology has the same respective
properties: Serre spectral sequences prove this for finite Postnikov
sections, which compute each fixed homology degree after a high
enough section.  The loop space of $Z[1/p]$ has
$\mathbb Z[1/p]$-module homology, and the loop space of
$Z_{\mathbb Q}$ has rational homology, by the same nilpotent
localization criterion.

Tensor and $\operatorname{Tor}_1$ vanish between a $p$-primary
group and a $\mathbb Z[1/p]$-module, and between a torsion group
and a rational vector space.  Integral K\"unneth and the simply
connected Whitehead theorem therefore give
\[
 G_p(Z)\wedge\Omega Z[1/p]\simeq *,\qquad
 F(Z)\wedge\Omega Z_{\mathbb Q}\simeq *.
 \tag{A34.5.4}
\]
The exact Ganea comparison in
\citet[Proposition 3.4, with Proposition 2.3 and Corollary 3.3]{A3ChataurScherer2006}
says that for $F\to E\to B$ the canonical map
$\operatorname{Cof}(F\to E)\to B$ has fiber
$F*\Omega B\simeq\Sigma(F\wedge\Omega B)$.
It follows from \textup{(A34.5.4)} that there are actual
cofiber sequences
\[
 G_p(Z)\longrightarrow Z\longrightarrow Z[1/p],\qquad
 F(Z)\longrightarrow Z\longrightarrow Z_{\mathbb Q}.
 \tag{A34.5.5}
\]
Combining them with \textup{(A34.5.3)} gives
\[
 \Sigma^\infty G_p(Z)\simeq\Gamma_p\Sigma^\infty Z,\qquad
 \Sigma^\infty F(Z)\simeq
       \Gamma_{\mathrm{tor}}\Sigma^\infty Z.
 \tag{A34.5.6}
\]

For spectra there is a natural equivalence
$\bigvee_p\Gamma_pE\simeq\Gamma_{\mathrm{tor}}E$.
To check it, these exact functors are smash products with their
values on the sphere.  The canonical map of those sphere values
is an equivalence after localization at every prime and after
rationalization: at $p$ all other summands vanish, and the
remaining fiber is $\operatorname{fib}(S_{(p)}\to H\mathbb Q)$.
Such localizations detect equivalences on every abelian homotopy
group, proving the assertion for all spectra.
The compatible maps $Z[1/p]\to Z_{\mathbb Q}$ give actual maps
$G_p(Z)\to F(Z)$.  Their wedge map has the preceding equivalence
as its suspension spectrum, by \textup{(A34.5.6)}.  Since its
source and target are simply connected, it is itself an equivalence:
\[
                       \bigvee_pG_p(Z)\simeq F(Z).
 \tag{A34.5.7}
\]
This uses the canonical fibers mapping into $Z$; no infinite product
of primary localizations or nonexistent reverse localization maps
are substituted for that wedge.

The exact membership test is now
\[
 X\in\mathcal A_W
 \quad\Longleftrightarrow\quad
 Z_{\mathbb Q}\in\mathcal A_W
 \ \text{and}\ G_p(Z)\in\mathcal A_W\text{ for every }p.
 \tag{A34.5.8}
\]
In the forward direction both $Z$ and $Z[1/p]$ are aspherical by
the telescope construction; their cofiber sequence makes $G_p(Z)$
aspherical.  In the reverse direction the wedge makes $F(Z)$
aspherical; its cofiber sequence with the rational term makes $Z$
aspherical.  Three reversible suspension reflections recover $X$.

\paragraph{Stable parts and the rational component.}
Each fundamental-localizer axiom and each reverse excision axiom
is preserved by intersection, proving the asserted properties of
$W_p$ and $W_{\mathbb Q}$.  Put $L=T(W)$.  Since $L$ is a
tensor ideal and the arithmetic functors are exact and preserve
colimits,
\[
 \Gamma_pL=L\cap I_p,\qquad
 (H\mathbb Q)\wedge L=L\cap\loc(H\mathbb Q).
 \tag{A34.5.9}
\]
For the first equality, $\Gamma_pE\in L\cap I_p$ whenever
$E\in L$, whereas $\Gamma_pE\simeq E$ for $E\in I_p$;
the rational equality is identical in form.

Generate $L$ by the suspension spectra of its aspherical spaces
$Y$, equivalently by the defining cofiber generators, using
reversibility.  For $Z=\Sigma^3Y$ the preceding construction puts
$G_p(Z)$ in $\mathcal A_{W_p}$, with suspension spectrum
$\Sigma^3\Gamma_p\Sigma^\infty Y$.
These generate $\Gamma_pL$, so $\Gamma_pL\subseteq T(W_p)$.
The opposite inclusion follows from $W_p\subseteq W$ and
$W_p\subseteq R(I_p)$.  Using $Z_{\mathbb Q}$ instead proves
the rational stable-part equality.

We use the checked pure-rational result REVLOC D.2.5 only with
$W_{\mathbb Q}$, whose whole stable part is rational.  Its crucial
argument is as follows.  A nonzero rational stable part supplies a
rational aspherical space with nonzero reduced homology.  Suspend
until the lowest nonzero homology degree $n\geq3$ is odd.
Rational Hurewicz and a nonzero rational linear functional on that
bottom homology split $S^n_{\mathbb Q}\simeq K(\mathbb Q,n)$
as a retract; this sphere identification is
\citet[Corollary 6.7.3]{A3MayPonto2012}.
Suspension reflection and rational cell attachments put every
connected rational space in the aspherical class, as in the full
proof of REVLOC D.2.5.  Thus $W_{\mathbb Q}=R(\loc(H\mathbb Q))$.
To cover an arbitrary pointed $X$ with rational suspension spectrum,
its triple suspension is simply connected and has rational integral
homology, hence is rational by the cited homology criterion; reversible
suspension reflection then recovers $X$ itself.  Thus no nilpotence
assumption on the original $X$ is hidden in this equality.
If its stable part is zero, it is $W_{\mathbb Z}$, since all
reversible fundamental localizers contain $R(0)=W_{\mathbb Z}$.

The join equality in \textup{(A34.5.1)} follows from
\textup{(A34.5.8)}: the rational term of an aspherical $Y$ lies
in $\mathcal A_{W_{\mathbb Q}}$, each primary term lies in
$\mathcal A_{W_p}$, and their cofiber reconstruction puts $Y$
in the join.  The converse inclusion is immediate.

\paragraph{Fixedness and a conditional categorical witness.}
If $W=R(L)$, then
$W_p=R(L)\cap R(I_p)=R(L\cap I_p)=R(T(W_p))$.
Conversely assume every $W_p$ is fixed, and take $X$ with
$\Sigma^\infty X\in L$.  For $Z=\Sigma^3X$,
\textup{(A34.5.6)} puts $\Sigma^\infty G_p(Z)$ in $T(W_p)$
and $\Sigma^\infty Z_{\mathbb Q}$ in $T(W_{\mathbb Q})$.
The assumed primary fixedness and proved rational fixedness make
all these spaces $W$-aspherical.  Equation \textup{(A34.5.8)}
does the same for $X$.  The reversible cofiber criterion gives
$R(L)\subseteq W$, proving \textup{(A34.5.2)}.

Explicitly, given $g\in R(T(W))\setminus W$, take $X=C(g)$
and $Z=\Sigma^3X$.  Its rational term is aspherical, so there
is a prime $p$ with $G_p(Z)\notin\mathcal A_W$.  Nevertheless
$\Sigma^\infty G_p(Z)\in T(W_p)$, and its reduced homology
is $p$-primary.  The category of singular simplices of $G_p(Z)$,
mapped to the terminal category, is therefore a conditional witness
in $R(T(W_p))\setminus W_p$.  If $T(W)$ has no nonzero finite
spectra, neither does $T(W_p)\subseteq T(W)$; its nonzero value
also follows from nonfixedness, since stable part zero forces
$W_p=W_{\mathbb Z}$.  This proves the reduction of A.3 for all
the localizers quantified there, without an accessibility premise.
\end{proof}


\begin{claim}{A34.6}{Proved}{A single accessible primary generator and an ordinary-null H-space test}\label{A346}
Fix Grothendieck universes $\mathcal U\in\mathcal V$, and interpret
spaces, spectra, categories and indexing sets as $\mathcal U$-small;
classes of such objects may be formed in $\mathcal V$.
The existential question A.3 is equivalent to the existence of the
following data:
\begin{enumerate}
\item a prime $p$ and a connected pointed CW homotopy type $V$ with
$p$-primary reduced integral homology;
\item the noncontractible infinite-loop space
$A=Q\Sigma^3V$, with
$L_A=\operatorname{loc}(\Sigma^\infty A)$ satisfying
$L_A\cap\operatorname{Ho}(\mathrm{Sp})^\omega=0$;
\item a noncontractible $d$-connected H-space $Z$, with
$p$-primary reduced integral homology, such that
\[
       P_AZ\simeq Z,
       \qquad \Sigma^\infty Z\in L_A.
       \tag{A34.6.1}
\]
\end{enumerate}
Here $n$ is the least degree with
$G=\widetilde H_n(A;\mathbb Z)\ne0$, and the finite integer $d$ is
specified by
\[
 d=\begin{cases}
       n-1,&G/pG\ne0,\\
       n,&G/pG=0.
    \end{cases}
 \tag{A34.6.2}
\]
One has $n\geq4$.  More precisely, for each prescribed finite
integer $m\geq1$, the same generator $A$ can be chosen with a
witness $Z_m$ satisfying \textup{(A34.6.1)} which is an
$m$-fold loop space.  The witness may depend on $m$; no single
infinite-loop witness is asserted.  For $m\geq2$, the unstable
algebra $H^*(Z_m;\mathbb F_p)$ cannot be finitely generated.
For these data the single-generator reversible
closure $W_A=\min\operatorname{Rev}(A\to *)$ is accessible,
$T(W_A)=L_A$, and the actual functor
\[
              \Delta/\operatorname{Sing}Z\longrightarrow e
              \tag{A34.6.3}
\]
belongs to $RT(W_A)\setminus W_A$.
More generally, for any connected primary $A$ which is
$(d-1)$-connected and kills every $d$-connected primary GEM,
\[
 H\in\mathcal A_{W_A}
 \quad\Longleftrightarrow\quad
 P_A\Sigma^dH\simeq *
 \quad\Longleftrightarrow\quad P_AH\simeq *
 \tag{A34.6.4}
\]
for every $d$-connected primary H-space $H$.
No countable model or uniform accessibility rank is asserted.
The data in \textup{(A34.6.1)} are still existential; this theorem
constructs them from an assumed counterexample, not unconditionally.
\end{claim}
\begin{proof}
\par\medskip\noindent\textbf{A small family of actual aspherical spaces.}\quad
Suppose A.3 has a witness.  By A34.5 it has a witness $W$ whose
stable part is contained in $I_p=\operatorname{loc}(\mathbb S/p)$
for one prime.  The cofiber criterion supplies a pointed CW type
$X$ with
\[
 \Sigma^\infty X\in T(W),\qquad X\notin\mathcal A_W.
 \tag{A34.6.5}
\]
It is connected: membership in $I_p$ makes its reduced integral
homology primary, while its reduced $H_0$ is free abelian.
Its suspension spectrum is nonzero, since $W_{\mathbb Z}\subseteq W$.

By A.2, the cofiber definition of $T(W)$ is equivalently the
localizing closure of suspension spectra of actual $W$-aspherical
pointed spaces.  We give the set argument on these spaces, avoiding
a subsequent choice of realizations of abstract spectra.  Let $\mathcal D$ consist
of all spectra lying in
$\operatorname{loc}\{\Sigma^\infty Y_i:i\in I\}$ for some
$\mathcal U$-small family of actual pointed CW models
$Y_i\in\mathcal A_W$.
The class $\mathcal D$ is closed under shifts and retracts.
Two witnesses can be replaced by their disjoint union, so it is
closed under cofibers of every morphism between its objects:
each generated subcategory is full.  For a $\mathcal U$-small
family of objects of $\mathcal D$, choose their witnessing families
in $\mathcal V$ and take their disjoint union.  That union is still
$\mathcal U$-small, by the closure axiom of a Grothendieck universe
for families indexed by one of its elements.  Thus $\mathcal D$
is also closed under all $\mathcal U$-small coproducts.
It contains every defining generator of $T(W)$ and is contained
in $T(W)$, proving $\mathcal D=T(W)$.
This is the small-subfamily argument already present in the supplied
REVLOC D.3.8; its use here begins with the actual aspherical models.

Choose such a family for $\Sigma^\infty X$, and put
\[
 V=\bigvee_{i\in I}Y_i,
 \qquad U=\Sigma^3V,
 \qquad A=QU.
 \tag{A34.6.6}
\]
All the $Y_i$ are connected and primary because
$W\subseteq R(T(W))\subseteq R(I_p)$.
They have $\mathcal U$-small pointed CW models, and the wedge,
suspension and a CW replacement of $QU$ remain $\mathcal U$-small.
Wedge and suspension closure give $U\in\mathcal A_W$.
By \citet[Proposition~3.1]{ABousfield1996}, a connected space $U$
ordinarily kills $QU$.  Ordinary $U$-acyclic spaces are
$W$-aspherical by the cone-attachment construction, as proved in
A34.1.  Hence $A\in\mathcal A_W$ and $W_A\subseteq W$.

The Snaith splitting, with its linear summand, gives
\[
 \Sigma^\infty QU\simeq
 \bigvee_{r\geq1}
      \bigl((\Sigma^\infty U)^{\wedge r}\bigr)_{h\Sigma_r}.
 \tag{A34.6.7}
\]
For the precise comparison, \citet[Theorems~2.7 and~6.1]{A3May1972}
identify the little-cubes construction $C_mU$ with
$\Omega^m\Sigma^mU$ for connected, well-pointed compactly generated
spaces, including these CW models without a countability bound.
Its $r$th filtration quotient is
$\mathcal C_m(r)_+\wedge_{\Sigma_r}U^{\wedge r}$.
The stable splitting of \citet[Theorem~1.1 and Section~2]{ASnaith}
is compatible with passage from $m$ to $m+1$ by
\citet[Remark~3.6]{ASnaith}.  At $m=\infty$, the space
$\mathcal C_\infty(r)$ is contractible and $\Sigma_r$-free
\citep[Corollary~4.6 and Theorem~4.8]{A3May1972}; hence the
quotient is the homotopy-orbit term in \textup{(A34.6.7)}.
Every localizing subcategory of spectra is a tensor ideal, so
all summands on the right lie in
$\operatorname{loc}(\Sigma^\infty U)$.
The $r=1$ summand gives the reverse inclusion.  Consequently
\[
 L_A=\operatorname{loc}(\Sigma^\infty U)
     =\operatorname{loc}\{\Sigma^\infty Y_i:i\in I\},
 \qquad \Sigma^\infty X\in L_A\subseteq T(W).
 \tag{A34.6.8}
\]
In particular $A$ is noncontractible and $L_A$ has no nonzero
finite spectra.  It is primary, and $A$ is $3$-connected;
these assertions follow from \textup{(A34.6.7)} and connective
stable Hurewicz, or from the corresponding connective spectrum
$\Sigma^\infty U$ and the nilpotent localization criterion.
The Galois relation gives $W_A\subseteq R(L_A)$, whereas its
aspherical generator gives $L_A\subseteq T(W_A)$.  Thus
\[
                         T(W_A)=L_A.
 \tag{A34.6.9}
\]
Because $W_A\subseteq W$, the space $X$ in
\textup{(A34.6.5)} is still outside its aspherical class.

\par\medskip\noindent\textbf{Accessibility of the reversible closure.}\quad
We recheck exactly the background conclusion being used.
For a pointed aspherical class $\mathcal C$ closed under pointed
homotopy colimits and all three cofiber implications, and an
uncountable regular $\kappa\in\mathcal U$, let
$\mathcal C^\kappa$ consist of the CW types for which every
subcomplex with fewer than $\kappa$ cells lies in a subcomplex
of that size belonging to $\mathcal C$.
The simultaneous enlargement proof in supplied REVLOC B.1--B.2
shows that $\mathcal C^\kappa\subseteq\mathcal C$, has the same
closure properties, and defines an accessible reversible cofiber
localizer.  Its fundamental generating set consists of all small
arrows whose cofibers are in $\mathcal C$, not merely of collapse
maps.  The proof enlarges a small subpair until its quotient is
an actual small $\mathcal C$-subcomplex, then uses the filtered
union of those subpairs.  Regularity and countable iteration keep
all sizes below $\kappa$; category-of-simplices representatives
still form a $\mathcal U$-small set.

Apply this to $\mathcal C=\mathcal A_{W_A}$, choosing $\kappa$
larger than the cell cardinality of $\Sigma A$.
The resulting sublocalizer is contained in $W_A$ and contains
$A\to *$, because its cofiber is $\Sigma A$.
Minimality among reversible localizers forces equality.
This is the full argument of supplied REVLOC B.4, rechecked here.
Thus $W_A$ is accessible.  The term ``single generator'' refers
to reversible closure; it does not assert that its underlying
fundamental localizer is generated by one collapse map.

\par\medskip\noindent\textbf{The precise finite suspension bound.}\quad
As $A$ is simply connected and noncontractible, it has a first
nonzero reduced integral homology group $G$ in some degree $n$;
otherwise the homological Whitehead theorem would make it
contractible.  Here $n\geq4$, $G$ is nonzero and primary, and
$A$ is $(n-1)$-connected by Hurewicz.
By \citet[Corollary~3.2(ii)]{ABousfield1996}, $A$ kills
$SP^\infty A$.  The Dold--Thom description makes $K(G,n)$ a
homotopy retract of this GEM, hence ordinarily $A$-acyclic.
If $G/pG\ne0$, \citet[Lemma~8.13]{ABousfield1997} says that
$K(G,n)$ kills every $(n-1)$-connected primary GEM.
If $G/pG=0$, that lemma additionally requires
$\pi_nY/p\pi_nY=0$; every $n$-connected GEM satisfies this
condition because $\pi_nY=0$.
Thus the value $d$ in \textup{(A34.6.2)} has both properties
needed in the more general assertion: $A$ is $(d-1)$-connected
and kills every $d$-connected primary GEM.  No divisible-bottom
case has been silently omitted.

\par\medskip\noindent\textbf{Desuspending the test on H-spaces.}\quad
We prove \textup{(A34.6.4)} under its general stated hypotheses.
A34.1 supplies the first equivalence.  Ordinary acyclicity is
closed under suspension, so it remains to prove
$P_A\Sigma^dH\simeq *\Rightarrow P_AH\simeq *$.
Put $Z=P_AH$.
The cone-attachment construction following
\citet[Theorem~1.1]{ABousfield1996} uses maps from nonnegative
suspensions of $A$.  Coning such maps preserves $d$-connectivity
of an input which has it: van Kampen gives simple connectivity,
the cofiber homology sequence preserves vanishing through degree
$d$, and Hurewicz then gives the asserted connectivity.
The same sequence preserves homology over $\mathbb Z[1/p]$.
These statements pass through the transfinite homotopy colimit,
so $Z$ is still $d$-connected and primary.
Nullification preserves finite products, as stated in the proof
of \citet[Theorem~1.7]{ABousfield1996}; hence $Z$ is an H-space.

The loops of an $A$-null space are $A$-null.  The universal
property of $H\to Z$ therefore gives
\[
 \operatorname{Map}_*(\Sigma^dZ,Z)
 \simeq\operatorname{Map}_*(Z,\Omega^dZ)
 \simeq\operatorname{Map}_*(H,\Omega^dZ)
 \simeq\operatorname{Map}_*(\Sigma^dH,Z)
 \simeq *.
 \tag{A34.6.10}
\]
Start with $j=d$, using \textup{(A34.6.10)}.  At a downward
induction step, let $1\leq j\leq d$ and assume only the already
established assertion
$\operatorname{Map}_*(\Sigma^jZ,Z)\simeq *$.  Apply \citet[Key Lemma~1.6]{ABousfield1996} with
$V'=\Sigma^{j-1}Z$ and target $Z$.
The source is connected, and $Z$ is a connected
$\Sigma V'$-null H-space, exactly as the lemma requires.
It gives
\[
 \operatorname{Map}_*(V',Z)
 \simeq\operatorname{Map}_*(SP^\infty V',Z)
 \simeq *.
 \tag{A34.6.11}
\]
The last equality holds because $SP^\infty V'$ is a
$d$-connected primary GEM, hence $A$-acyclic, whereas $Z$
is $A$-null.  Descending to $j=1$ makes
$\operatorname{Map}_*(Z,Z)$ contractible.  Its identity is
nullhomotopic, so $Z\simeq *$.  This proves the whole equivalence
including its mapping-space components.  It does not commute
nullification with suspension.

\par\medskip\noindent\textbf{Producing the H-space witness from the assumed witness.}\quad
Return to $X$ in \textup{(A34.6.5)}, and let
\[
                         H=\Omega\Sigma^{d+1}X.
 \tag{A34.6.12}
\]
It is a $d$-connected H-space.  The reduced-product filtration,
or the case $n=1$ of \citet[Theorem~1.1]{ASnaith} together
with \citet[Corollary~6.2]{A3May1972}, has stable layers $(\Sigma^d\Sigma^\infty X)^{\wedge r}$.
Thus \textup{(A34.6.8)} and tensor-ideal closure give
$\Sigma^\infty H\in L_A$, and $H$ is primary.
The original James formulation is
\citet[Section~1 and Theorem~5.6]{A4James1955}; no countability
hypothesis is needed here because the cited Snaith theorem
applies to every connected compactly generated CW model.
The suspension--loop adjunction triangle exhibits
$\Sigma^{d+1}X$ as a homotopy retract of $\Sigma H$.
If $H$ were $W_A$-aspherical, suspension, retract closure and
reversible desuspension would put $X$ in that class, a
contradiction.  Equation \textup{(A34.6.4)} therefore makes
$Z=P_AH$ noncontractible.

The previous paragraph of the proof establishes its required
connectivity, primary homology and H-space structure.
The cone-attachment construction starts from $H$ and attaches
cofibers built from suspensions of $A$.  Its suspension spectrum
therefore remains in $L_A$ at every stage and at the transfinite
homotopy colimit.  Consequently this $Z$ satisfies
\textup{(A34.6.1)}.

For the finite-loop strengthening, fix $m\geq1$ and use instead
\[
 H_m=\Omega^m\Sigma^{m+d}X,\qquad Z_m=P_AH_m.
 \tag{A34.6.13}
\]
May's approximation and the finite-$m$ Snaith splitting have
layers
$\mathcal C_m(r)_+\wedge_{\Sigma_r}(\Sigma^dX)^{\wedge r}$.
The coefficient action is free, so these are homotopy-orbit
constructions; tensor-ideal and homotopy-colimit closure put
all their suspension spectra in $L_A$.
The adjunction triangle now makes $\Sigma^{m+d}X$ a retract of
$\Sigma^mH_m$.  Exactly the preceding argument therefore makes
$Z_m$ noncontractible, $d$-connected, primary, $A$-null, and
stably in $L_A$.
For its actual loop structure, \citet[Theorem~5.2]{ABousfield1997}
states that $L_f\Omega Y\simeq\Omega L_{\Sigma f}Y$ for every
pointed map $f$ and connected space $Y$.
Taking $f=A\to *$ and iterating, with all intervening spaces
connected here, gives
\[
 Z_m\simeq\Omega^m P_{\Sigma^m A}\Sigma^{m+d}X.
\]
This proves the strengthening separately for every finite $m$,
without passing to a limit of the witnesses.

\par\medskip\noindent\textbf{The finite-generation exclusion in this normal form.}\quad
More generally, let $A,d$ satisfy the hypotheses of
\textup{(A34.6.4)}, with $d\geq1$, and let a $d$-connected primary
$A$-null space $Z$ be the loop space of a connected H-space $U$.
Suppose that $H^*(Z;\mathbb F_p)$ is finitely generated as an
unstable algebra.  Put $Y=P_{BC_p}Z$.  Theorem~5.2 of
\citet{ABousfield1997} gives
\[
             Y\simeq\Omega P_{\Sigma BC_p}U.
\]
Nullification preserves finite products, so the delooping on the
right is again a connected H-space.  The cone construction of
$P_{BC_p}$ preserves simple connectivity and reduced
$\mathbb Z[1/p]$-acyclicity; hence $Y$ is simply connected and
primary.  By \citet[Theorem~7.2]{A4CCSDeconstruct2007}, its total
mod-$p$ cohomology is finite dimensional.
\citet[Theorems~1.1 and~2.2]{A3Hemmi1991} prove, for every prime,
that a loop space on an H-space is a quasi-$C_\infty$ space and
that a simply connected quasi-$C_p$ space with finite mod-$p$
cohomology has zero reduced mod-$p$ cohomology.  Thus
$\widetilde H^*(Y;\mathbb F_p)=0$.
The homological universal coefficient sequence surjects
$H_{q+1}(Y;\mathbb F_p)$ onto
$\widetilde H_q(Y;\mathbb Z)[p]$.
Every nonzero $p$-primary torsion group has an element of order
$p$, including in the divisible case.  Consequently $Y$ is
integrally acyclic and, being simply connected, contractible.

Lemma~7.1 and Theorem~5.5 of
\citet{A4CCSDeconstruct2007} identify the actual fiber of
$Z\to P_{BC_p}Z$ as a $p$-torsion H-Postnikov piece with only
finitely many nonzero homotopy groups: the lemma supplies the
finite Krull stage and the Lannes-$T$ finite-type hypotheses
from finite unstable-algebra generation.
Since the base $Y$ is contractible, this fiber is $Z$ itself.
Every nonzero Eilenberg--Mac Lane stage is $d$-connected and
primary, so is killed by $A$.  Finite fiber-extension closure
of ordinary $A$-acyclicity gives $P_AZ\simeq *$; $A$-nullity
then gives $Z\simeq *$.  For each $m\geq2$, the witness $Z_m$
above is a loop space on a connected H-space and is
noncontractible, proving the asserted exclusion.  This argument
uses neither sphere-nullity nor a completion of $Z$.

Conversely, suppose the stated data are given.  Equation
\textup{(A34.6.4)} shows that a noncontractible $A$-null $Z$
cannot be $W_A$-aspherical.  Its suspension spectrum lies in
$T(W_A)=L_A$, so it is $RT(W_A)$-aspherical.
The category-of-simplices comparison realizes this as the actual
functor \textup{(A34.6.3)}.  Its stable part has no nonzero finite
spectra and its reversible localizer is accessible, proving A.3
under the assumed data.  This completes both directions.
\end{proof}

\paragraph{Background and scope of the new reduction.}
The small-subfamily argument and accessibility of small reversible
closures are supplied REVLOC D.3.8 and B.4, respectively; they are
reproved above rather than claimed as new standalone results.
The new conclusion is their combination with the universal prime
reduction and the ordinary-null H-space test for the full
existential question A.3.  The proof gives no bound on the size
of $V$, $A$ or $Z$ and no uniform accessibility rank.  It therefore
does not settle the fixed-stable-part set-size question A.4.
The original existence question remains open, and no priority
claim is made.
